\section{Methodology}
\label{sec:methodology}

\subsection{The loop}
% DECISIONS.md s7 (cold-start), s18 (retention/budget mode), s6 (one weight channel), s8 (seeds/determinism)
Each round: the adversary generates a batch of flows and a per-batch cost record (flows,
packets, bytes, flow-time); the auto-labeler stamps every row in the batch (Section~\ref{sec:threat-model});
a defense hook, when one is installed, sets a per-row \texttt{sample\_weight} and may quarantine
the whole batch; the batch is added to the accumulated training set; the model is
\textbf{retrained from scratch} on seed data plus every admitted batch to date; the fresh model
is scored on \texttt{trusted\_eval}. Retraining cold-start each round, rather than warm-starting
from the previous round's model, keeps rounds comparable and prevents poisoning dynamics from
being confounded with optimization path dependence.
% DECISIONS.md s18: retention=accumulate (default, worst case for the defender: honeypot
% influence never expires), budget_mode=fixed_ratio (the per-round batch size is set so the
% honeypot's share of the training set reaches a target ratio by the final round)
Two knobs govern how honeypot data accumulates. \emph{Retention} is \texttt{accumulate} by
default: every admitted batch is kept for the rest of the run, which is both the design the
honeypot-retraining literature proposes ("continuous adaptation") and the worst case for the
defender, since influence only grows. \emph{Budget mode} is \texttt{fixed\_ratio}: the
per-round batch size is set so the honeypot's share of the training set reaches a target ratio
by the final round, so every point on a damage curve has a known, comparable poison share.
% DECISIONS.md s6: single sample_weight channel; no class_weight="balanced", no scale_pos_weight,
% no oversampling/SMOTE -- A1 shifts the class balance by design, and automatic rebalancing would
% respond to that shift and confound the measured effect
No model in this codebase uses automatic class rebalancing (\texttt{class\_weight="balanced"},
XGBoost's \texttt{scale\_pos\_weight}, oversampling, SMOTE). A1 works by injecting
malicious-labeled volume, which shifts the class balance; an automatic rebalancer would respond
to that shift and the measured attack effect would be entangled with a rebalancing artifact.
\texttt{sample\_weight} is the only channel by which any row's influence on the fit is adjusted,
and every defense in Section~\ref{sec:defenses} rides on it.
% DECISIONS.md s8: seed_everything (Python, NumPy, PyTorch), deterministic kernels, float32
% features; the model seed varies by round (seed*1000+round) but is shared across arms at a
% given (seed, round), so paired same-seed deltas across arms isolate the treatment
Every fit seeds Python, NumPy and PyTorch and forces deterministic kernels. The model seed
varies by round but is shared across every arm at a given (seed, round), so a paired,
same-seed delta between an arm and the control isolates the treatment rather than retraining
noise.

\subsection{Datasets and partitions}
% DECISIONS.md s1: within_day_temporal splits within each (day, family) block, not each whole
% day, because a whole-day split places every attack family entirely inside one partition and
% destroys A4 as a variable. Fractions: benign 30/20/30/20, attack 40/-/40/20 (no prod_benign
% for attack). Four eval arms: seen_both / seed_only / honeypot_only / novel.
Data comes from CICIDS2017 (\texttt{MachineLearningCVE}) and a synthetic generator that ships
with the codebase, both partitioned identically into four disjoint sets: \texttt{seed\_train}
(the defender's trusted starting data), \texttt{prod\_benign} (unused here), \texttt{honeypot\_pool}
(what the adversary can draw from and what genuine decoy traffic would supply), and
\texttt{trusted\_eval} (held out, never written to by the loop). The split is within each
(day, family) block, not each whole day: benign rows go 30/20/30/20 to
seed\_train/prod\_benign/honeypot\_pool/trusted\_eval, attack rows 40/40/20 to
seed\_train/honeypot\_pool/trusted\_eval. This gives every attack family a presence in every
partition (\texttt{seen\_both}), which a whole-day split cannot, since CICIDS2017 runs each
family on its own day.
% PAPER_PLAN.md writing rule + DECISIONS.md s14 erratum: MachineLearningCVE has no Timestamp
% column; clean_flows synthesizes one at one row per second in file order. Every "temporal"
% split, guard-(b) buffer and burst boundary on CICIDS in this paper is therefore a file-order
% cut, not a wall-clock one.
\textbf{The CICIDS split is not temporal.} The \texttt{MachineLearningCVE} CSVs carry no
timestamp column; a timestamp is synthesized at one row per second in file order. Every cut,
buffer and burst boundary computed on CICIDS in this paper is therefore a file-order cut, not
a wall-clock one. Section~\ref{sec:degeneracy} shows this distinction does not matter to the
result: a uniformly shuffled re-split reproduces the same feature-space distances.
% DECISIONS.md s3: seed_only must never be merged with novel -- only seed_only is a controlled
% manipulation (A4), and collapsing it hides A4 entirely.
Four eval-family arms are tracked, never folded into an aggregate: \texttt{seen\_both} (in both
seed\_train and honeypot\_pool, the literature's regime), \texttt{seed\_only} (A4's controlled
arm), \texttt{honeypot\_only} (the loop's most interesting arm: a family the detector saw only
via the decoy), and \texttt{novel} (in neither).

% DECISIONS.md s24(b) + synthetic.py docstring: the generator emits raw-CICIDS-shaped data
% with well-separated class geometry (so a clean S0 loop can demonstrably improve) and injects
% the same corruption classes the real CSVs carry (NaN/Inf rates, negative durations, duplicate
% rows) so the cleaning pipeline is exercised end to end. It has no model of attacker effort or
% cost -- every synthetic flow's duration/packets/bytes come from the same class-conditional
% distribution regardless of scenario, which is why Section~\ref{sec:defenses}'s cost-based
% defense cannot separate benign from attack rows on this dataset.
The synthetic generator emits data in the same shape as the real CSVs, with well-separated
class geometry so a clean loop can demonstrably improve, and injects the same corruption
classes (non-finite rates, negative durations, duplicate rows) so the cleaning pipeline is
exercised the same way on both datasets. It has no model of attacker cost: a flow's duration,
packet count and byte count are drawn from the same class-conditional distribution regardless
of which scenario produced it.

\begin{table}[t]
  \centering
  \caption{Partition sizes (rows), near-duplicate grid 0.005. Attack rows have no
  \texttt{prod\_benign} share.}
  % Source: results/phase0/cicids/partition_summary.csv, results/phase0/partition_summary.csv
  \label{tab:partitions}
  \small
  \begin{tabular}{lrrrr}
    \toprule
    Partition & \multicolumn{2}{c}{CICIDS2017} & \multicolumn{2}{c}{Synthetic} \\
    & benign & attack & benign & attack \\
    \midrule
    seed\_train      & 344{,}823 & 78{,}877  & 20{,}383 & 8{,}241 \\
    prod\_benign     & 229{,}139 & --        & 13{,}098 & --      \\
    honeypot\_pool   & 344{,}221 & 130{,}166 & 19{,}951 & 7{,}993 \\
    trusted\_eval    & 229{,}695 & 52{,}054  & 13{,}262 & 4{,}535 \\
    \bottomrule
  \end{tabular}
\end{table}

\subsection{Leakage guards}
% DECISIONS.md s2: three guards on within_day_temporal -- (a) global near-duplicate removal
% (two half-offset grid snaps, grid 0.005 -- s11, s15), (b) a temporal guard band around
% internal cut times, (c) nearest-neighbour distance from every trusted_eval row to its
% nearest honeypot_pool row of the same class, reported as a percentile distribution for BOTH
% classes but gating leak_warning on the ATTACK class only.
Three guards run on every partition build. (a) Near-duplicate rows are removed globally, before
splitting, at grid resolution 0.005 (grid sweep and rationale in
Appendix~\ref{sec:appendix}). (b) Rows within a buffer of an internal cut are dropped
(burst-level splitting, a stronger version of this guard, is also in the appendix; it does not
change Section~\ref{sec:degeneracy}'s result). (c) For every \texttt{trusted\_eval} row, the
nearest-neighbour distance in normalized feature space to its nearest \texttt{honeypot\_pool}
row of the same class is computed and reported for both classes — this is the metric behind
every "realized distance" and "twin fraction" number in this paper.
% DECISIONS.md s2, s15: leak_warning gating rationale moved to Appendix for space; guard (c)'s
% definition above is kept in the main body since S5/S6 depend on readers knowing what it is.
Why the resulting \texttt{leak\_warning} is gated on the attack class alone, and what that does
and does not imply about the benign side, is in Appendix~\ref{sec:appendix}.

\subsection{Threshold policy}
% DECISIONS.md s5: calibrate to a target FPR (0.01) on a benign-stratified validation split
% carved from training, conservative under ties. Two modes, both reported: fixed (round-0
% threshold held) shows A1 damage as rising FPR; recalibrated (every round) shows it as
% collapsing TPR.
No model uses a 0.5 cutoff. Each model's threshold is calibrated to a target FPR (0.01) on a
validation split carved from its own training data. Two threshold modes are reported for every
run, because A1's damage surfaces in different places under each: \emph{fixed} calibrates once
at round 0 and holds that threshold, so damage appears as rising FPR on trusted benign;
\emph{recalibrated} recalibrates every round on trusted benign, so FPR stays capped by
construction and damage appears as collapsing TPR instead.
% DECISIONS.md s19.1: the round-0 validation split is full of near-copies of training rows on
% CICIDS (s16), so naive calibration is optimistic and the threshold too low. Fix: calibrate
% only on validation rows with no same-class training twin within tau=0.1 (per-feature RMS,
% normalized space). This discards 94% of benign validation rows and moves control fixed-FPR
% from 0.046 to 0.007 (RF).
On CICIDS2017 this calibration is not honest by default: the validation split is carved from
training data that, by Section~\ref{sec:degeneracy}'s measurement, is full of near-copies of
what it is meant to hold out, so a threshold calibrated against it is optimistic and lands too
low (control fixed-threshold FPR 0.046 against a 0.01 target). We calibrate instead only on
validation rows with no same-class training twin within a distance $\tau=0.1$ (the same
per-feature RMS metric as guard (c)). This discards 94\% of benign validation rows and moves
the control's fixed-threshold FPR to 0.007. Every CICIDS number in this paper after
Section~\ref{sec:degeneracy} uses this calibration.

\subsection{Arms}
% DECISIONS.md s23 table + s19.2: control (no ingestion, retrain on seed_train each round only,
% model seed varies by round so it measures real retrain noise); S0 (genuine attack rows,
% correctly labeled); A1 (benign rows, jittered, stamped malicious); plus three matched null
% controls that each remove one candidate mechanism: a1truth (the same rows as A1, true label:
% removes the label conflict), s0j (attack rows with A1's jitter, stamped malicious: same
% volume and prior shift, no label conflict), junk (marginal-shuffled bulk, stamped malicious:
% volume alone, no structure).
Every experiment compares an arm against a same-seed \emph{control} (no honeypot ingestion,
retrain on seed\_train alone each round; the model seed still varies by round, so the control
measures real retrain-to-retrain noise, not zero). \textbf{S0} (clean) feeds genuine attack
rows, correctly labeled. \textbf{A1} (benign mimicry) feeds benign rows, jittered by a
controlled amount, stamped malicious by the auto-labeler. Section~\ref{sec:attack} isolates
A1's mechanism with three matched null controls, each removing exactly one candidate cause:
\texttt{a1truth} (A1's own rows, correctly labeled — removes the label conflict),
\texttt{s0j} (attack rows carrying A1's jitter, stamped malicious — same volume and class-prior
shift, no label conflict), and \texttt{junk} (marginal-shuffled bulk, stamped malicious — volume
with no structure at all).

\subsection{Models, grids and metrics}
% DECISIONS.md s24 "Grid": RF and XGBoost hyperparameters (FAST_HYPERPARAMS in
% src/dloop/loop/rounds.py); 20 rounds, 5 seeds throughout.
Two models are trained per round: a random forest (25 trees, max depth 16) and XGBoost (60
trees), both at library defaults otherwise. Every result in this paper is 20 rounds, 5 seeds,
unless stated otherwise.
% src/dloop/loop/data.py LoopDataConfig defaults: seed_train_rows=30000,
% eval_benign_rows=40000, eval_attack_cap_per_family=4000 -- the fixed subsample the loop runs
% on each round (cheap enough to retrain from scratch every round, thousands of times over a
% sweep), drawn once per dataset and shared by every arm/seed/round so seed-to-seed variance
% reflects model and adversary randomness, not a changing evaluation set.
Each round retrains on a fixed subsample of the partitions in Table~\ref{tab:partitions}: up to
30{,}000 \texttt{seed\_train} rows, evaluated against up to 40{,}000 trusted benign rows and up
to 4{,}000 \texttt{trusted\_eval} rows per attack family, drawn once per dataset (not per seed)
so between-seed variance reflects model and adversary randomness rather than a moving
evaluation target.
% DECISIONS.md s4: A1 poison ratio grid is log-spaced, 0/0.005/0.01/0.02/0.05/0.1/0.2/0.5,
% because if A1 does real damage it shows at the low end and that is where resolution matters.
% Section 6's damage-vs-distance sweep uses ratios {0.05,0.1,0.2,0.5} and jitters
% {0,0.01,0.03,0.1,0.3,0.7,1.5}; s18: jitter is roughly log-spaced in REALIZED distance, and
% realized distance (not the jitter knob) is the reported x-axis.
The A1 poison-ratio grid is log-spaced (finer at the low end, where damage first appears if it
appears at all): $\{0, 0.005, 0.01, 0.02, 0.05, 0.1, 0.2, 0.5\}$ generally; the
damage-versus-distance sweep in Section~\ref{sec:attack} uses ratios
$\{0.05, 0.1, 0.2, 0.5\}$ and jitter $\{0, 0.01, 0.03, 0.1, 0.3, 0.7, 1.5\}$, roughly log-spaced
in \emph{realized} distance.
% DECISIONS.md s18 "Realized fidelity, not the knob": A1's x-axis is the median NN distance
% from the injected poison to the FULL trusted_eval benign set (guard-(c) metric, normalized
% space, per-feature RMS), computed per batch/round from the whole realized-distance
% distribution (p0-p100 recorded, median used for the x-axis); asserted disjoint from
% trusted_eval by row content.
The jitter knob is a generation parameter, not what is reported: every A1 plot's x-axis is the
\emph{realized} mimicry distance — the median nearest-neighbour distance (guard-(c) metric,
per-feature RMS, normalized space) from the injected poison rows to the full trusted benign
set, measured directly, not assumed from the jitter setting.
% DECISIONS.md s24 "Axes": R = 1 - mean(dFPR_defended)/mean(dFPR_undefended), A1 raw fidelity
% (jitter 0), ratio 0.5, fixed threshold. U = mean(gain_defended)/mean(gain_undefended), S0 on
% synthetic, honeypot_only TPR gain over control, ratios 0.05 and 0.2, fixed threshold.
Section~\ref{sec:defenses} reports two summary quantities against these arms. \emph{Recovery},
$R = 1 - \overline{\Delta\text{FPR}}_{\text{defended}} / \overline{\Delta\text{FPR}}_{\text{undefended}}$,
is measured on A1 at raw benign fidelity (jitter 0), poison ratio 0.5, fixed threshold — the
setting where undefended damage is largest. \emph{Retention},
$U = \overline{\text{gain}}_{\text{defended}} / \overline{\text{gain}}_{\text{undefended}}$, is
the honest loop's own \texttt{honeypot\_only} TPR gain over control on S0 (synthetic, ratios
0.05 and 0.2, fixed threshold), defended over undefended: how much of the honest loop's benefit
a defense leaves intact.
% DECISIONS.md s8: paired same-seed deltas at the final round; every table/figure delta in this
% paper is arm_metric[seed] - control_metric[seed], same seed, same round, then mean/sd over
% the 5 seeds and a paired t = mean / (sd / sqrt(5)).
Every reported delta is paired by seed: for each of the 5 seeds, an arm's final-round metric
minus the control's final-round metric at that same seed, then the mean, standard deviation and
paired $t = \bar{d} / (s_d/\sqrt{5})$ over the 5 differences.
